Cell Painting microscopy captures how cells change after a chemical or genetic perturbation.
Connecting these images to the perturbations that produced them could make large imaging screens
easier to search and interpret, but the task remains difficult because biological effects are
subtle and technical variation is substantial.
We introduce MorphoCLIP, a contrastive model that links Cell Painting profiles with text
descriptions of compounds, CRISPR knockouts and ORF overexpressions. The model keeps its vision
and language backbones frozen and trains only a compact cross-channel module and projection
layers, so it can be trained on a single consumer GPU.
On held-out CPJUMP1 data, MorphoCLIP searches in both directions: from a cell image to its
perturbation description and from a description to matching cell images. In both cases, a correct
match appears among the top ten results much more often than expected by chance. Adding a
replicate-alignment loss makes profiles from repeated experiments more consistent, although this
improvement does not yet translate into reliable gene--compound matching. Gene-aware labels and
plate correction also show no consistent retrieval benefit.
These findings suggest that text supervision can help organize chemical and genetic Cell Painting
data. Matching compounds with genetic perturbations, however, remains an open problem.
